\section{Mathematical Preliminaries}
\label{sec:preliminaries}

In this section, we review the mathematical background necessary for understanding stereographic encoding, including the stereographic projection, the Bloch sphere representation of qubits, and Möbius transformations.

\subsection{Stereographic Projection}

The stereographic projection is a mapping from the extended complex plane $\C \cup \{\infty\}$ (the Riemann sphere) to the unit sphere $\Sphere \subset \R^3$.

\begin{definition}[Stereographic Projection]
Let $\Sphere = \{(u,v,w) \in \R^3 : u^2 + v^2 + w^2 = 1\}$ be the unit sphere. The stereographic projection $\pi: \C \to \Sphere \setminus \{(0,0,1)\}$ is defined by
\begin{align}
    u &= \frac{2x}{|z|^2 + 1}, \\
    v &= \frac{2y}{|z|^2 + 1}, \\
    w &= \frac{|z|^2 - 1}{|z|^2 + 1},
\end{align}
where $z = x + iy \in \C$.
\end{definition}

The inverse stereographic projection $\pi^{-1}: \Sphere \setminus \{(0,0,1)\} \to \C$ is given by
\begin{equation}
    x = \frac{u}{1-w}, \quad y = \frac{v}{1-w}.
\end{equation}

This projection is conformal (angle-preserving) and bijective, mapping circles in $\C$ to circles on $\Sphere$. The point at infinity in $\C$ corresponds to the north pole $(0,0,1)$ of the sphere.

\subsection{The Bloch Sphere}

A general single-qubit pure state can be written as
\begin{equation}
    \ket{\psi} = \psi_0 \ket{0} + \psi_1 \ket{1},
\end{equation}
where $\psi_0, \psi_1 \in \C$ satisfy the normalization condition $|\psi_0|^2 + |\psi_1|^2 = 1$.

\begin{definition}[Bloch Sphere Representation]
Any single-qubit pure state can be parameterized as
\begin{equation}
    \ket{\psi(\theta, \varphi)} = \cos\left(\frac{\theta}{2}\right)\ket{0} + e^{i\varphi}\sin\left(\frac{\theta}{2}\right)\ket{1},
\end{equation}
where $\theta \in [0,\pi]$ is the polar angle and $\varphi \in [0, 2\pi)$ is the azimuthal angle. This parameterization maps pure states bijectively to points on the unit sphere $\Sphere$ (the Bloch sphere).
\end{definition}

The density matrix of a pure state is $\rho = \ket{\psi}\bra{\psi}$, which can be expanded in the Pauli basis:
\begin{equation}
    \rho = \frac{1}{2}(I + \vec{r} \cdot \vec{\sigma}),
\end{equation}
where $\vec{r} = (\langle \PauliX \rangle, \langle \PauliY \rangle, \langle \PauliZ \rangle)$ is the Bloch vector with $|\vec{r}| = 1$ for pure states, and $\vec{\sigma} = (\PauliX, \PauliY, \PauliZ)$ are the Pauli matrices:
\begin{equation}
    \PauliX = \begin{pmatrix} 0 & 1 \\ 1 & 0 \end{pmatrix}, \quad
    \PauliY = \begin{pmatrix} 0 & -i \\ i & 0 \end{pmatrix}, \quad
    \PauliZ = \begin{pmatrix} 1 & 0 \\ 0 & -1 \end{pmatrix}.
\end{equation}

\subsection{Möbius Transformations}

\begin{definition}[Möbius Transformation]
A Möbius (or fractional linear) transformation is a function $f: \C \cup \{\infty\} \to \C \cup \{\infty\}$ of the form
\begin{equation}
    f(z) = \frac{az + b}{cz + d},
\end{equation}
where $a,b,c,d \in \C$ satisfy $ad - bc \neq 0$.
\end{definition}

Möbius transformations have several important properties:
\begin{itemize}
    \item They are bijective and form a group under composition.
    \item They are conformal (angle-preserving).
    \item They are automorphisms of the Riemann sphere.
    \item They map circles and lines to circles and lines.
    \item The inverse of $f(z) = (az+b)/(cz+d)$ is $f^{-1}(w) = (-dw+b)/(cw-a)$.
\end{itemize}

There is a natural correspondence between Möbius transformations and $2 \times 2$ matrices. The transformation $f(z) = (az+b)/(cz+d)$ can be associated with the matrix
\begin{equation}
    M = \begin{pmatrix} a & b \\ c & d \end{pmatrix}.
\end{equation}
Composition of Möbius transformations corresponds to matrix multiplication (up to scalar multiples).

\subsection{Chebyshev Polynomials}

The Chebyshev polynomials of the first and second kind play a central role in our analysis of quantum signal processing.

\begin{definition}[Chebyshev Polynomials]
The Chebyshev polynomials of the first kind $T_n(x)$ and second kind $U_n(x)$ are defined by
\begin{align}
    T_n(\cos\theta) &= \cos(n\theta), \\
    U_n(\cos\theta) &= \frac{\sin((n+1)\theta)}{\sin\theta}.
\end{align}
Equivalently, for $x \in [-1,1]$:
\begin{align}
    T_n(x) &= \cos(n \arccos x), \\
    U_n(x) &= \frac{\sin((n+1)\arccos x)}{\sin(\arccos x)} = \frac{\sin((n+1)\arccos x)}{\sqrt{1-x^2}}.
\end{align}
\end{definition}

These polynomials satisfy the recurrence relations:
\begin{align}
    T_0(x) &= 1, \quad T_1(x) = x, \quad T_{n+1}(x) = 2xT_n(x) - T_{n-1}(x), \\
    U_0(x) &= 1, \quad U_1(x) = 2x, \quad U_{n+1}(x) = 2xU_n(x) - U_{n-1}(x).
\end{align}
